\appendix

\section{Supplementary Legal, Economic, and Methodological Notes}
\label{sec_appendix}

\subsection{Compensatory measures in the Portuguese setting}

The Ericeira source document emphasized a practical difficulty that deserves to remain visible in the final paper: although compensatory measures are central in Portuguese environmental and maritime negotiations, the legal framework does not provide a simple operational definition for surf-related losses. Decree-Law No.\ 140/99 frames compensatory measures as actions required to preserve the coherence of the Natura 2000 network when impacts cannot be fully avoided. In practice, this means that compensation is not merely a cash-transfer issue. It must be connected to the ecological and territorial effects of the project, the habitats affected, and the populations or uses that depend on those habitats.

For the surf case, this creates an evidentiary challenge. If the relevant loss is linked to changes in wave quality, submerged morphology, event viability, or broader destination attractiveness, then the affected side must translate those impacts into arguments that can be recognized in licensing and negotiation arenas. This is one reason why the article treats institutional representation as part of the strategic problem itself rather than as a background detail.

\subsection{Additional notes on Ericeira and surf-tourism value}

The Ericeira material also compiled evidence that surfing generates a wider local economic web than a narrow sports classification would suggest. The source highlights the expansion of restaurants, bars, surf camps, surf houses, and specialized accommodation connected to the area's international surf reputation. It also reports expenditure evidence indicating that surfers spend not only on accommodation, but also on food, beverages, and transport, which supports the interpretation that wave quality has multiplier effects across the local service economy.

For this reason, the surf side in the model should not be read as a small recreational niche. It is better understood as a coalition around a place-based economic asset. This framing is consistent with the article's use of Peniche as an order-of-magnitude benchmark and strengthens the interpretation that expected losses may be rationally perceived as broad local-economic losses rather than only sporting inconvenience.

\subsection{Additional questions raised by the Peniche note}

The file of questions stored in the older \texttt{artigo} folder usefully sharpened some issues that are only implicit in the main text. First, it framed the dispute as a tension between two public-policy objectives rather than a simple conflict between a private developer and local recreation. On one side sits renewable-energy expansion, justified by decarbonization and the broader transition toward a high share of renewable generation. On the other sits the local economic policy logic tied to surf-driven tourism, place promotion, and community livelihood.

Second, the note stressed that environmental assessment cannot be satisfied by studying an ``average wave'' if the relevant economic resource is a specific surfable wave profile. For surfing, the decisive question is not only whether waves continue to exist, but whether the altered wave remains suitable for the same type of practice, the same user profile, and the same events. A change in wave type may preserve water motion while still destroying the local surf value that anchors tourism demand.

Third, the note sharpened an institutional question that fits directly into the model: if compensation becomes part of the negotiation, who is the legitimate receiver? The surfing field may contain local associations, event structures, informal service networks, schools, accommodation operators, and broader tourism actors whose economic dependence on surfing is real but organizationally fragmented. This reinforces the article's argument that institutional alignment between surf and tourism is not secondary. It may be a precondition for negotiating, valuing loss, and receiving compensation in a coherent way.

For these reasons, the old note was not treated as a separate article draft. Its useful content was absorbed here as a set of analytical questions that strengthen the framing of loss, assessment, and representation in the main paper.

\subsection{Extended game-theory background}

The support review contained a broader methodological discussion that goes beyond what was necessary in the main body of the paper. Three points are especially useful as a supplement.

First, the conflict studied here is related to, but not identical with, the classic tragedy-of-the-commons structure. In commons problems, individually rational choices may degrade a shared resource even when all actors would benefit from restraint \citep{hardin1968tragedy,ostrom_2015}. The present paper studies a more specific confrontation in which one coalition advances a project and another resists it. Even so, the policy intuition is similar: without institutional design, decentralized strategic behavior can produce an inferior collective outcome.

Second, the review stressed the importance of distinguishing Nash equilibrium from Pareto superiority. A cell can be stable because no player wants to deviate unilaterally and still be socially worse than another available outcome. This distinction is the analytical core of the present article, because the baseline worksheet makes conflict stable even though negotiated acceptance yields a higher total payoff.

Third, the review argued that policy intervention should not be restricted to punitive logic. In game-theory terms, strategic outcomes can be shifted by reducing the payoff from non-cooperation, by increasing the payoff from cooperation, or by changing the cost structure that sustains conflict. Applied to offshore wind and surf-tourism, this means that relocation, redesign, technical mitigation, procedural mediation, subsidies for adaptation, and compensation architecture are all legitimate candidates for changing the game.

\subsection{Modelling caveats and future extensions}

The longer methodological notes also raised several caveats that remain valid. The current formulation is static, aggregates heterogeneous actors into two coalitions, and treats the dispute as a one-shot decision. Real maritime conflicts are more complex. They may unfold sequentially, repeat across licensing stages, and include reputation effects, coalition fragmentation, and learning over time.

There is also room to model uncertainty more richly. Instead of a single probability of offshore success, future work may use intervals, scenarios, fuzzy assessments, or evidence-weighted probabilities. Another extension is to compare litigation-cost regimes more explicitly, including the American rule, the British rule, and asymmetric cost-shifting arrangements. The support material also suggested that repeated interactions could introduce strategies based on reciprocity, memory, and reputation, which would be especially relevant in long coastal-governance processes.

Finally, the review noted that positive incentives may be as important as sanctions. In policy terms, it may be more effective to create credible gains from negotiated coexistence than to rely exclusively on penalties for resistance or non-cooperation. This observation is consistent with the results of the exploratory analysis delivered with the paper.

\subsection{Immediate analytical tasks suggested during model construction}

The project notes also identified a short list of analytical tasks that should guide the next iteration of the article. The first is to characterize more explicitly the conditions under which the local side should fight and the conditions under which acceptance becomes rational. Although the current paper already provides threshold inequalities, a more systematic presentation of the fighting and non-fighting regions would improve interpretability.

The second task is to compare global loss more directly across solution concepts. In practical terms, this means quantifying the distance between the Nash equilibrium, the Pareto-best cell, and any broader notion of aggregate welfare used by the study. For the current worksheet, this can be expressed as the difference in total payoff between the conflict equilibrium and the negotiated outcome, but the note correctly suggests extending this to the losses borne by each side separately.

The third task is graphical. The notes propose plotting the model in terms of the probability of offshore success ($P$), varying the local loss ($L$) and compensation ($G$), and comparing cases in which offshore returns are one, two, or three orders of magnitude larger than local surf-tourism income. This is a useful next step because it would convert the current threshold logic into a visual sensitivity map of when resistance is individually rational and when compensation ceases to be effective.
